\documentclass[10pt,a4paper]{amsart}
\usepackage[margin=19mm]{geometry}
\usepackage{amsmath,amssymb,mathtools}
\usepackage[scr=rsfs,cal=euler]{mathalfa}
\usepackage{xfrac}
\usepackage[unicode,hidelinks,bookmarks=false]{hyperref}
\newcommand{\ie}{i.e.\ }
\newcommand{\cf}{cf.\ }
\newcommand{\resp}{respectively\ }
\newcommand{\Subsection}{\subsection}
\providecommand{\nobreakdash}{\nobreak\hbox{-}}
\setlength{\emergencystretch}{2em}
\multlinegap0pt
\usepackage{bm}
\usepackage{bbm}
\usepackage{dsfont}
\usepackage{xspace}
\usepackage{cancel}
\usepackage{mathtools}
\usepackage{amsthm}
\makeatletter
\@ifundefined{theorem}{\newtheorem{theorem}{Theorem}}{}
\@ifundefined{proposition}{\newtheorem{proposition}[theorem]{Proposition}}{}
\makeatother
\def\N{\mathbb{N}}
\def\R{{\mathbb {R}}}
\newcommand{\abs}[1]{\lvert#1\rvert}
\title[Local and global solvability of the Grushin heat equation wi]{Local and global solvability of the Grushin heat equation with mixed nonlinear memory and reaction terms}
\author{Ahmad Z. Fino, Arlucio Viana}
\date{Source: arXiv:2507.13547v1 (CC BY 4.0)}
\begin{document}
\maketitle

\noindent\emph{Source and licence.} Local and global solvability of the Grushin heat equation with mixed nonlinear memory and reaction terms. Version arXiv:2507.13547v1 (CC BY 4.0), distributed under the Creative Commons Attribution 4.0 International licence, as recorded in the arXiv licence metadata for this exact version and shown on the abstract page https://arxiv.org/abs/2507.13547v1. Full rights notice: licenses/arxiv-2507.13547-CC-BY-4.0.txt.\par\smallskip

\noindent\textit{Editorial scope.} Counted unit: the Theorem whose source label is \texttt{comparison} in the article (source0.tex lines 626--634, source label \texttt{comparison}), delivered together with the article's own complete original proof of it (source0.tex lines 636--675, one \texttt{proof} environment of 40 lines). No mathematics is written, repaired or re-derived: statement and proof are the article's own text line for line. Required context retained uncounted: * (lines 146--154); properties (lines 198--214); solop (lines 452--454); growth1 (lines 616--618). Every definition, display and label of those lines is retained unchanged. Omitted from this package and not redistributed: 1--145, 155--197, 215--451, 455--615, 619--625, 635--635, 676--970. Nothing inside a retained excerpt is omitted or altered: every retained body is delivered exactly as the article's lines are written, so no omission receipt of any kind accompanies this package. Citations: the retained excerpts carry no citation command at all, so no bibliography entry of the article is transcribed and no reference is redistributed. Reference adaptation: none was needed and none was made; every reference inside the delivered unit resolves inside it, so no unresolved reference or citation is produced. Printed slips kept literal: no printed word, symbol, hypothesis or number of the counted statement or of its proof was altered. Layout and class: the article's own class is \texttt{article} with packages including amsmath, amsfonts, latexsym, amscd, amssymb, epsfig, amsthm, enumerate, tikz, color, hyperref, natbib; the wrapper is the lane's pinned \texttt{amsart} editorial wrapper at 10pt on A4, which restates the theorem-like environments the retained text needs and adds the article's own macro definitions that the retained text needs (\texttt{\textbackslash N}, \texttt{\textbackslash R}, \texttt{\textbackslash abs}), with extra wrapper packages (bm, bbm, dsfont, xspace, cancel, mathtools). The delivered numbering is the unit's own. Assets: no retained span contains a figure environment, a graphics-inclusion command, a PDF-page inclusion, a table or a TikZ picture; no image file is copied into this package, the package declares no asset dependency and no image file is redistributed with it. Licence: version arXiv:2507.13547v1 of this article is distributed under the Creative Commons Attribution 4.0 International licence, as recorded in the arXiv licence metadata for this exact version and shown on the abstract page https://arxiv.org/abs/2507.13547v1. A full rights notice is delivered at licenses/arxiv-2507.13547-CC-BY-4.0.txt. Characters: every retained excerpt is pure ASCII, so the delivered engine, XeLaTeX with its default Latin Modern text font, reports no missing character.
\begin{equation}\label{*}
	\left\{\begin{array}{ll}
		\displaystyle {u_{t}-\Delta_{\mathcal{G}}
			u = k_1 \int_0^t(t-s)^{-\gamma}\abs{u}^{p_1-1}u(s)\,\mathrm{d}s} + k_2|u|^{p_2-1}u, &\displaystyle {t>0,z\in {\mathbb{R}}^{N+k},}\\
		{}\\
		\displaystyle{u(x,0)=  u_0(z),\qquad\qquad}&\displaystyle{z=(x,y)\in {\mathbb{R}}^{N+k},}
	\end{array}
	\right.
\end{equation}

\begin{proposition}\label{properties}
	There is a unique semigroup $(S_{\mathcal{G}}(t))_{t\geq0}$ associated to $\Delta_{\mathcal{G}}$ and satisfying the following properties:
	\begin{itemize}
		\item[$(i)$]  For all $\varphi\in L^p(\R^{N+k})$, $1\leq p\leq \infty$, we have
		$$S_{\mathcal{G}}(t)\varphi(x,y)=\int_{\R^{N+k}} K(x,w,y-z;t) \varphi(w,z)\,dw\,dz,\qquad (x,y)\in \R^{N+k},\,\,t>0.$$
		\item[$(ii)$] $(S_\mathcal{G}(t))_{t\geq0}$ is a contraction semigroup on $L^p(\mathbb{R}^{N+k})$, $1\leq p\leq \infty$, which is strongly continuous for $1\leq p<\infty$. 
		\item[$(iii)$] $(S_\mathcal{G}(t))_{t\geq0}$  is a strongly continuous semigroup on $C_0(\mathbb{R}^{N+k})$.
		\item[$(iv)$] For every $v\in X$, where $X$ is either $L^p(\mathbb{R}^{N+k})$ for $1\leq p< \infty$ or $C_0(\mathbb{R}^{N+k})$, the map $t\longmapsto S_\mathcal{G}(t)v$ is continuous from $[0,\infty)$ into $X$.
		\item[$(v)$] $\displaystyle\int_{\mathbb{R}^{N+k}}K(x,0,y;t)  \,dx\,dy=1$, for all $t>0$.
		\item[$(vi)$] The heat kernel satisfies
		$$K(x,x_0,y;t)>0,\qquad \textrm{for all}\,\, (x, x_0,y;t)\in\R^{2N+k}\times(0,\infty).$$
		Hence, $S_\mathcal{G}(t)\varphi\geq 0$ whenever $\varphi\geq0$, for all $t\geq0$.
		\item[$(vii)$] $K(rx,rx_0,r^2y;r^2t)=r^{-(N+2k)}K(x,x_0,y;t)$, for all $(x, x_0,y;t)\in\R^{2N+k}\times(0,\infty)$, $r>0$.
		\item[$(viii)$]  For all $(x, x_0,y;t)\in\R^{2N+k}\times(0,\infty)$, we have
		$$K(x,x_0,-y;t)=K(x,x_0,y;t)\quad\mathrm{and}\quad K(x,x_0,y;t)=K(x_0,x,y;t).$$
	\end{itemize}
\end{proposition} 

	\begin{equation}\label{solop}
		\Lambda(u)(t):=S_{\mathcal{G}}(t)u_0 +k_1 \int_{0}^{t}S_{\mathcal{G}}(t-s) \int_0^s(s-\tau)^{-\gamma}\abs{u(\tau)}^{p_1-1}u(\tau) \, \mathrm{d}\tau \,\mathrm{d}s + k_2\int_{0}^{t}S_{\mathcal{G}}(t-s) |u(s)|^{p_2-1}u(s) \mathrm{d}s.
	\end{equation}

\begin{equation}\label{growth1}
	|f(u)-f(v)| \leq c|u-v|(|u|^{p_1-1} + |v|^{p_1-1}),\quad\quad |g(u)-g(v)| \leq c|u-v|(|u|^{p_2-1} + |v|^{p_2-1}).
\end{equation}

\begin{theorem}\label{comparison}
	Let $u_0,v_0\in X$, $k_1,k_2\geq 0$, and $f,g, \tilde{f}, \tilde{g}$ functions satisfying \eqref{growth1}.
	\begin{itemize}
		\item[(i)] If $0\leq u_0\leq v_0$, and $f(\mu)\leq \tilde{f}(\nu)$, $g(\mu)\leq \tilde{g}(\nu)$, for all $0\leq\mu\leq \nu$, then 
		$$0\leq u(t;u_0,f,g) \leq v(t;v_0,\tilde{f},\tilde{g})\quad\text{for all}\,\, t\in[0,\tilde{T}],$$
		where $\tilde{T}>0$ is a common time of existence for $u$ and $v$.
		\item[(ii)]  If $v$ is a mild supersolution of \eqref{*}, and $u$ is a mild solution with $u_0\geq0$, then $0\leq u\leq v$ for all $t\in[0,T]$.
	\end{itemize}
\end{theorem}

\begin{proof} We begin by defining the Picard sequences:
	$$u_{1}(t) =S_\mathcal{G}(t)u_0,\qquad u_{k}(t) =\Lambda_{fg} (u_{k-1})(t),\,\,\, k\geq2 ,$$
	$$v_{1}(t) =S_\mathcal{G}(t)v_0 ,\qquad v_{k}(t) =\widetilde{\Lambda}_{\tilde{f}\tilde{g}} (v_{k-1})(t),\,\,\, k\geq2 , $$
	where $\Lambda_{fg}$ and $\widetilde{\Lambda}_{\tilde{f}\tilde{g}}$ refer to the operator introduced in equation \eqref{solop}, emphasizing the dependence on the nonlinearities in the right-hand side:
	\begin{equation}\label{lambdafg}
		\Lambda_{fg}(u)(t):=S_{\mathcal{G}}(t)u_0 +k_1 \int_{0}^{t}S_{\mathcal{G}}(t-s) \int_0^s(s-\tau)^{-\gamma}f(u(\tau))\, \mathrm{d}\tau \,\mathrm{d}s +k_2 \int_{0}^{t}S_{\mathcal{G}}(t-s)g(u(s)) \mathrm{d}s,
	\end{equation}
	$$\widetilde{\Lambda}_{\tilde{f}\tilde{g}}(u)(t):=S_{\mathcal{G}}(t)v_0 +k_1 \int_{0}^{t}S_{\mathcal{G}}(t-s) \int_0^s(s-\tau)^{-\gamma}\tilde{f}(u(\tau))\, \mathrm{d}\tau \,\mathrm{d}s +k_2 \int_{0}^{t}S_{\mathcal{G}}(t-s)\tilde{g}(u(s)) \mathrm{d}s.$$
	By Proposition \ref{properties}-(vi), we know that the semigroup $S_\mathcal{G}(t)$ preserves order:
	\begin{equation}\label{monotonicity}
		S_\mathcal{G}( t)\varphi \leq S_\mathcal{G}(t)\psi,\quad \hbox{whenever}\,\,\, \varphi\leq \psi.
	\end{equation} 
	Hence, $0\leq u_1(t)\leq v_1(t)$. We now apply induction to show that $0\leq u_k(t)\leq v_k(t)$, for all $k\in\N$. Assume this holds for some fixed $k$. Then, using the monotonicity of the semigroup $S_\mathcal{G}(t)$ (cf. \eqref{monotonicity}) together with the structural conditions imposed on the nonlinearities $f, \tilde{f},g,\tilde{g}$, we obtain
	$$0\leq u_{k+1}(t)=\Lambda_{fg}(u_k)(t)\leq \widetilde{\Lambda}_{\tilde{f}\tilde{g}}(v_k)(t)=v_{k+1}(t).$$
	Hence, the inequality is preserved at each iteration step. This proves part $(i)$, noting that the corresponding mild solutions are obtained as pointwise limits of the Picard sequences, depending on the functional setting:
	\begin{itemize}
		\item  If $X=C_0(\mathbb{R}^{N+k})$, then the convergence is uniform on compact sets and, in particular, pointwise. Thus,
		$$u(t;u_0,f,g) = \lim\limits_{n\rightarrow\infty} u_n(t),\qquad  v(t;v_0,\tilde{f},\tilde{g}) = \lim\limits_{n\rightarrow\infty} v_n(t),$$
		where the limits are taken pointwise.
		\item  If $X=L^q(\mathbb{R}^{N+k})$,  then by norm convergence and standard results, there exists a common subsequence along which both sequences converge pointwise almost everywhere. That is,
		$$u(t;u_0,f,g) = \lim\limits_{k\rightarrow\infty} u_{n_k}(t),\qquad  v(t;v_0,\tilde{f},\tilde{g}) = \lim\limits_{k\rightarrow\infty} v_{n_k}(t),$$ 
		where the limits are taken almost everywhere along the same subsequence $(n_k)_k$.\\
	\end{itemize}
	To prove part $(ii)$, it suffices to observe that since $v$ is a mild supersolution, we can construct a non-increasing Picard sequence $(w_k)_k$ defined by
	$$
	w_{1}(t) =v(t),\qquad\quad w_{k}(t) =\Lambda_{fg} (w_{k-1})(t),\,\,\hbox{for}\,\, k\geq2.
	$$
	By part $(i)$, we have
	$$w_{2}(t) =\Lambda_{fg} (w_1)(t)=\Lambda_{fg} (v)(t)\leq v(t)=w_1(t),$$ 
	and an induction argument using the monotonicity property of $\Lambda_{fg}$ (as established in part $(i)$) yields that the sequence $(w_k)_k$  is non-increasing. Moreover, since the nonlinearities $f,g$ are nonnegative and the semigroup $S_\mathcal{G}(t)$ preserves nonnegativity, each $w_k(t)$ remains bounded, satisfying $0\leq w_k(t)\leq v(t)$ for all $k\geq1$.
	\begin{itemize}
		\item  If the functional space is $X=C_0(\mathbb{R}^{N+k})$, then $(w_k)_k$ converges uniform on compact sets and, in particular,
		pointwise to a limit function $w(t)$. By standard monotone convergence results in Banach spaces, the convergence is strong in $C([0,T],C_0(\mathbb{R}^{N+k})))$. Passing to the limit in the iteration yields
		$$w(t)= \Lambda_{fg} (w)(t)\leq v(t),$$
		in particular, $w$ is a mild solution of problem \eqref{*}. By uniqueness, we conclude that $u=w\leq v$.
		\item  If instead $X=L^q(\mathbb{R}^{N+k})$,  then the sequence $(w_k)_k$ converges in norm in $L^q$, and in particular, there exists a subsequence $(w_{k_j})_j$ that converges almost everywhere to the same limit function $w$. Passing to the limit almost everywhere in the integral formulation gives
		$$w(t)= \Lambda_{fg} (w)(t)\leq v(t),$$
		so again $w$ is a mild solution of \eqref{*}. Uniqueness implies $u=w\leq v$.
	\end{itemize} 
\end{proof}

\end{document}
